\documentclass[12pt,a4paper]{article}
\usepackage[utf8]{inputenc}
\usepackage[vietnamese]{babel}
\usepackage{graphicx}
\usepackage{geometry}
\usepackage{hyperref}
\usepackage{listings}
\geometry{a4paper, margin=1in}

\title{\textbf{BÁO CÁO RBL MÔN CƠ SỞ DỮ LIỆU (DBI202)}\\
\large TIẾN ĐỘ THỰC HIỆN TỪ TUẦN 5 ĐẾN TUẦN 7\\
\small Đề tài: Quản lý Ca làm việc và Chấm công (Staff Shift Management)}
\author{Nhóm Thực Hiện: Nhóm 01}
\date{\today}

\begin{document}

\maketitle

\tableofcontents
\newpage

\section{Báo cáo Tuần 5: Chọn đề tài \& Xác định vai trò của DB}

\subsection{Tên đề tài nghiên cứu}
\textbf{Hệ thống Quản lý Ca làm việc và Chấm công nhân viên (Staff Shift Management)}

\subsection{Lý do chọn đề tài}
Hiện nay, các cửa hàng bán lẻ, chuỗi F\&B hay doanh nghiệp vừa và nhỏ sử dụng rất nhiều nhân viên làm việc theo ca (part-time và full-time). Việc chấm công bằng Excel thủ công hoặc chép tay vào sổ rất dễ gây ra nhầm lẫn giờ làm, tính sai tiền lương, nhất là với các ca làm đêm vắt sang ngày hôm sau.

Nhóm chọn đề tài này nhằm mục đích thiết kế một cơ sở dữ liệu quan hệ chặt chẽ giúp tự động hóa việc lưu trữ lịch làm việc, ghi nhận thẻ chấm công hằng ngày và hỗ trợ tính toán thù lao một cách chính xác.

\subsection{Xác định các thực thể và mối quan hệ}
Mô hình dữ liệu của hệ thống gồm 3 thực thể chính:
\begin{itemize}
    \item \textbf{STAFF (Nhân viên):} Quản lý mã nhân viên (\texttt{UserID}), họ tên, vai trò và mức lương thù lao theo giờ (\texttt{HourlyRate}).
    \item \textbf{SHIFT (Ca làm việc):} Quản lý các khung ca định sẵn (\texttt{ShiftID}, \texttt{ShiftName}, \texttt{Time\_Range}).
    \item \textbf{TIMECARD (Thẻ chấm công):} Lưu nhật ký quẹt thẻ hằng ngày gồm ngày làm (\texttt{WorkDate}), giờ vào (\texttt{CheckIn}), giờ ra (\texttt{CheckOut}).
\end{itemize}

\noindent \textbf{Các mối quan hệ giữa thực thể:}
\begin{itemize}
    \item Quan hệ giữa \texttt{STAFF} và \texttt{SHIFT} là quan hệ Nhiều - Nhiều (M-N). Một nhân viên đăng ký làm nhiều ca, và một ca có nhiều nhân viên. Nhóm dùng bảng trung gian \textbf{SCHEDULES} để kết nối.
    \item Quan hệ giữa \texttt{STAFF} và \texttt{TIMECARD} là quan hệ Một - Nhiều (1-N).
    \item Quan hệ giữa \texttt{SHIFT} và \texttt{TIMECARD} là quan hệ Một - Nhiều (1-N).
\end{itemize}

\subsection{Câu hỏi nghiên cứu (Research Question - RQ)}
\textit{“Làm thế nào để tổ chức cơ sở dữ liệu quan hệ giúp quản lý phân ca nhân viên chặt chẽ, đồng thời hỗ trợ khai thác dữ liệu để tự động hóa việc tính toán tổng số giờ làm việc thực tế và xác định nhân viên đi muộn hay đúng giờ?”}

\subsection{Mô tả sơ bộ nguồn dữ liệu}
\begin{itemize}
    \item \textbf{Dữ liệu tĩnh:} Danh sách nhân sự và các khung giờ ca làm việc do bộ phận quản lý nhập vào CSDL.
    \item \textbf{Dữ liệu động:} Lịch sử ra/vào làm việc phát sinh hằng ngày từ thẻ chấm công của nhân viên.
\end{itemize}

\section{Báo cáo Tuần 6: Công cụ \& Quy trình triển khai Dataset}

\subsection{Giới thiệu công cụ sử dụng}
Nhóm thống nhất sử dụng các công cụ và kỹ thuật sau:
\begin{itemize}
    \item \textbf{Hệ quản trị CSDL:} Microsoft SQL Server (dùng công cụ SSMS để viết câu lệnh SQL khởi tạo bảng, khóa chính, khóa ngoại và ràng buộc kiểm tra).
    \item \textbf{Ngôn ngữ lập trình xử lý:} JavaScript (chạy trên môi trường Node.js để viết script đọc dữ liệu thô và xử lý toán học về thời gian).
    \item \textbf{Trình soạn thảo code:} Visual Studio Code.
    \item \textbf{Nền tảng báo cáo:} Overleaf (LaTeX) và Prism.openchat.ai.
\end{itemize}

\subsection{Sơ đồ luồng công việc (Workflow)}
Luồng thu thập và xử lý tập dữ liệu của hệ thống diễn ra qua 4 bước tuần tự:
\begin{enumerate}
    \item \textbf{Khởi tạo CSDL:} Tạo các bảng và thiết lập ràng buộc toàn vẹn dữ liệu trên SQL Server (ví dụ constraint \texttt{CHECK (HourlyRate > 0)}).
    \item \textbf{Thu thập dữ liệu thô:} Ghi nhận các bản ghi quẹt thẻ ra vào ca của nhân viên vào bảng \texttt{TIMECARD}.
    \item \textbf{Làm sạch và tính toán ban đầu:} Script Node.js (\texttt{process\_timecards.js}) truy vấn bảng chấm công, thực hiện tính số giờ làm (\texttt{TotalHours}) và kiểm tra giờ vào làm để gán nhãn \texttt{On Time} hoặc \texttt{Late}.
    \item \textbf{Trình bày đầu ra:} Dữ liệu sau xử lý xuất ra file chuẩn và kết nối lên giao diện Web để quản lý theo dõi.
\end{enumerate}

\section{Báo cáo Tuần 7: Thu thập dữ liệu thực tế \& Xử lý ban đầu}

\subsection{Quá trình tạo tập dữ liệu thô (Raw Data)}
Nhóm sử dụng câu lệnh \texttt{INSERT INTO} tạo tập dữ liệu thử nghiệm (Mock Data) sát với thực tế, bao gồm:
\begin{itemize}
    \item Nhân viên đi làm đúng giờ (Check-in lúc \texttt{07:55:00} ca sáng).
    \item Nhân viên đi làm muộn (Check-in lúc \texttt{08:15:00} ca sáng).
    \item Bản ghi làm ca đêm vắt qua ngày hôm sau (Check-in lúc \texttt{22:00:00} và Check-out lúc \texttt{06:00:00} sáng hôm sau).
\end{itemize}

\subsection{Giải thích mã nguồn làm sạch dữ liệu (\texttt{process\_timecards.js})}

\subsubsection{Logic tính tổng số giờ làm (\texttt{calculateTotalHours})}
Dữ liệu giờ vào/ra được tách bằng hàm \texttt{.split(':')} để quy đổi toàn bộ sang đơn vị phút. 

\noindent \textbf{Khó khăn thực tế sinh viên gặp phải:} Khi nhóm test thử bản ghi làm ca đêm từ 22h đêm đến 6h sáng hôm sau, công thức lấy \texttt{(giờ ra - giờ vào)} cho kết quả bị âm \texttt{(6 - 22 = -16 tiếng)}. Sau khi debug kiểm tra lại, nhóm đã sửa mã nguồn bằng cách thêm điều kiện kiểm tra: nếu tổng số phút nhỏ hơn 0 thì cộng thêm \texttt{24 * 60} phút (tức bù thêm một ngày). Nhờ đó kết quả tính toán ra chuẩn 8.00 tiếng.

\subsubsection{Logic phán đoán trạng thái đi làm (\texttt{evaluateStatus})}
Hệ thống dùng lệnh điều kiện \texttt{If-Else}: Đối với Ca sáng (quy định bắt đầu lúc 08:00), nếu giờ Check-in lớn hơn 8 hoặc đúng 8 giờ nhưng số phút lớn hơn 0 thì gán nhãn \texttt{'Late'} (Đi muộn), ngược lại gán nhãn \texttt{'On Time'}.

\subsection{Kết quả đạt được sau Tuần 7}
Nhóm đã thu được tập dữ liệu sạch hoàn chỉnh, đã loại bỏ lỗi tính toán thời gian ca đêm, lưu dưới dạng file \texttt{processed\_timecards.csv} sẵn sàng cho các bước tối ưu ở Tuần 8.

\section{Kết luận}
Tính đến Tuần 7, nhóm đã hoàn thành đúng tiến độ các mục tiêu đề ra trong hướng dẫn RBL môn học DBI202. Việc kết hợp giữa CSDL quan hệ SQL Server và script Node.js xử lý ngoại lệ thời gian giúp hệ thống vận hành ổn định, sát với yêu cầu thực tế của doanh nghiệp.

\end{document}
